% GENERATED FILE. Edit canonical lessons, then run npm run generate:books.
% canonical-source-hash: fnv1a64:48999076a73c7a5f
% canonical-lessons: PA-C180-lok, PA-C180-us, PA-C180-thaan, PA-C180-kamre, PA-C180-vadde, PA-C180-place-review, PA-C180-payoff

\chapter{People in the Big Room}
\label{ch:pa-a2-place-reference}
\hlchaptermodality{180}

\begin{chapteropening}
\textbf{By the end of this chapter:} \emph{I can understand and give one simple practical-place note about people in a room.}

Retrieve a one-sentence place note through listening, speaking, reading, and bounded Gurmukhi writing after an atom-free review.
\end{chapteropening}

\section[\texorpdfstring{\pa{ਲੋਕ} (lok)}{ਲੋਕ (lok)}]{\pa{ਲੋਕ} — people}

\label{lesson:PA-C180-lok}



\begin{quote}
\emph{Hear:} \emph{lok}; YOU SAY: \textbf{\pa{ਲੋਕ}} Start by ear. This one word means \emph{people}.
\end{quote}

\begin{grammarlens}[title={Grammar Lens: one group word}]
\begin{quote}
\textbf{\pa{ਲੋਕ}} — \emph{lok} — people
\end{quote}

Keep this as one useful word. No new plural rule is needed today.
\end{grammarlens}

\begin{cousinweb}
\textbf{\pa{ਲੋਕ}} (\emph{lok}) is Sanskrit \emph{loka}, \textquotedblleft{}world, people\textquotedblright{}, taken over as a learned loan. Mahan Kosh identifies the Sanskrit word; CDIAL 11119 gives the same meanings.
\end{cousinweb}

\subsection*{Your turn}
\begin{itemize}
\raggedright
  \item \emph{Hear:} \emph{lok}; YOU SAY: \textbf{\pa{ਲੋਕ}}
  \item \emph{Read it:} \textbf{\pa{ਲੋਕ}}
\end{itemize}

\begin{tcolorbox}[breakable,colback=teal!4,colframe=teal!35!black,title={Before you move on}]
Say \emph{people}. (\textbf{\pa{ਲੋਕ}}.)

Sources: Punjabi University, \href{https://dic.learnpunjabi.org/default.aspx?look=\%E0\%A8\%B2\%E0\%A9\%8B\%E0\%A8\%95}{RCPLT dictionary}, lists \textbf{\pa{ਲੋਕ}} as people or public. \href{https://punjabi.com/dictionary/wordmeaning?script=roman\&search=\%E0\%A8\%B2\%E0\%A9\%8B\%E0\%A8\%95}{Mahan Kosh} identifies Sanskrit \emph{loka}; Turner gives its senses in \href{https://dsal.uchicago.edu/dictionaries/soas/}{CDIAL 11119}.
\end{tcolorbox}
\section[\texorpdfstring{\pa{ਉਸ} (us)}{ਉਸ (us)}]{\pa{ਉਸ} — pointing to that place}

\label{lesson:PA-C180-us}



\begin{quote}
Say \emph{people} once. (\textbf{\pa{ਲੋਕ}}.)
\end{quote}

\begin{grammarlens}[title={Grammar Lens: one pointing word}]
\begin{quote}
\textbf{\pa{ਉਸ} \pa{ਬਾਗ਼} \pa{ਵਿੱਚ} \pa{ਲੋਕ} \pa{ਹਨ।}} — \emph{us bāġ vic lok han} — people are in that garden.
\end{quote}

Here \textbf{\pa{ਉਸ}} points to \emph{that} garden. Treat the whole line as one place model; it does not teach every form of \emph{that}.
\end{grammarlens}

\subsection*{Your turn}
\begin{itemize}
\raggedright
  \item \emph{Hear:} the garden model; YOU SAY: \textbf{\pa{ਉਸ} \pa{ਬਾਗ਼} \pa{ਵਿੱਚ} \pa{ਲੋਕ} \pa{ਹਨ।}}
  \item \emph{Read it:} \textbf{\pa{ਉਸ} \pa{ਬਾਗ਼} \pa{ਵਿੱਚ} \pa{ਲੋਕ} \pa{ਹਨ।}}
\end{itemize}

\begin{tcolorbox}[breakable,colback=teal!4,colframe=teal!35!black,title={Before you move on}]
Say \emph{that garden}, then \emph{people}. (\textbf{\pa{ਉਸ} \pa{ਬਾਗ਼}}, \textbf{\pa{ਲੋਕ}}.)

Source: Punjabi University, \href{https://dic.learnpunjabi.org/default.aspx?look=\%E0\%A8\%89\%E0\%A8\%B8}{RCPLT dictionary}, lists \textbf{\pa{ਉਸ}} as \emph{that} in a conjunctive form.
\end{tcolorbox}
\section[\texorpdfstring{\pa{ਥਾਂ} (thaan)}{ਥਾਂ (thaan)}]{\pa{ਥਾਂ} — one place word}

\label{lesson:PA-C180-thaan}



\begin{quote}
Retrieve \textbf{\pa{ਉਸ}} and \textbf{\pa{ਲੋਕ}} by ear.
\end{quote}

\begin{grammarlens}[title={Grammar Lens: one location noun}]
\begin{quote}
\textbf{\pa{ਥਾਂ}} — \emph{thā̃} — place
\end{quote}

\begin{quote}
\textbf{\pa{ਉਸ} \pa{ਥਾਂ}} — \emph{us thā̃} — that place
\end{quote}

Keep this as a two-word pointing phrase, not a new rule for every location noun.
\end{grammarlens}

\begin{cousinweb}
\textbf{\pa{ਥਾਂ}} (\emph{thā̃}) is a tadbhava, inherited and reshaped by sound change. It comes through Prakrit \emph{thāma} from Sanskrit \emph{sthāman}, \textquotedblleft{}station, place\textquotedblright{} (CDIAL 13760).
\end{cousinweb}

\subsection*{Your turn}
\begin{itemize}
\raggedright
  \item \emph{Hear:} \emph{us thā̃}; YOU SAY: \textbf{\pa{ਉਸ} \pa{ਥਾਂ}}
  \item \emph{Read it:} \textbf{\pa{ਉਸ} \pa{ਥਾਂ}}
  \item \emph{Say it:} \textbf{\pa{ਲੋਕ}}
\end{itemize}

\begin{tcolorbox}[breakable,colback=teal!4,colframe=teal!35!black,title={Before you move on}]
Say \emph{that place}. (\textbf{\pa{ਉਸ} \pa{ਥਾਂ}}.)

Sources: Punjabi University, \href{https://dic.learnpunjabi.org/default.aspx?look=\%E0\%A8\%A5\%E0\%A8\%BE\%E0\%A8\%82}{RCPLT dictionary}, gives \textbf{\pa{ਥਾਂ}} as place or location. Turner derives it through Prakrit \emph{thāma} from Sanskrit \emph{sthāman} in \href{https://dsal.uchicago.edu/dictionaries/soas/}{CDIAL 13760}.
\end{tcolorbox}
\section[\texorpdfstring{\pa{ਕਮਰੇ} (kamre)}{ਕਮਰੇ (kamre)}]{\pa{ਕਮਰੇ} — room in one location phrase}

\label{lesson:PA-C180-kamre}



\begin{quote}
Say \textbf{\pa{ਉਸ} \pa{ਥਾਂ}}, then the old room word \textbf{\pa{ਕਮਰਾ}}.
\end{quote}

\begin{grammarlens}[title={Grammar Lens: one room phrase}]
\begin{quote}
\textbf{\pa{ਕਮਰਾ}} — \emph{kamrā} — a room (the old word)
\end{quote}

\begin{quote}
\textbf{\pa{ਕਮਰੇ} \pa{ਵਿੱਚ}} — \emph{kamre vic} — in the room
\end{quote}

The ending changes inside this one location phrase. Use this model first; do not infer every possible room ending from it.
\end{grammarlens}

\subsection*{Your turn}
\begin{itemize}
\raggedright
  \item \emph{Hear:} \emph{kamre vic}; YOU SAY: \textbf{\pa{ਕਮਰੇ} \pa{ਵਿੱਚ}}
  \item \emph{Read it:} \textbf{\pa{ਕਮਰੇ} \pa{ਵਿੱਚ}}
  \item \emph{Say it:} \textbf{\pa{ਉਸ} \pa{ਥਾਂ}}
\end{itemize}

\begin{tcolorbox}[breakable,colback=teal!4,colframe=teal!35!black,title={Before you move on}]
Say \emph{in the room}, then \emph{that place}. (\textbf{\pa{ਕਮਰੇ} \pa{ਵਿੱਚ}}, \textbf{\pa{ਉਸ} \pa{ਥਾਂ}}.)

Source: Punjabi University, \href{https://punjabiuniversity.ac.in/Pages/Department.aspx?dsenc=129}{library description}, uses \textbf{\pa{ਕਮਰੇ} \pa{ਵਿੱਚ}} in its phrase for 'in the study room'.
\end{tcolorbox}
\section[\texorpdfstring{\pa{ਵੱਡੇ} (vadde)}{ਵੱਡੇ (vadde)}]{\pa{ਵੱਡੇ} — big in one room phrase}

\label{lesson:PA-C180-vadde}



\begin{quote}
Say the old \textbf{\pa{ਵੱਡਾ}}, then \textbf{\pa{ਕਮਰੇ} \pa{ਵਿੱਚ}}.
\end{quote}

\begin{grammarlens}[title={Grammar Lens: one agreeing phrase}]
\begin{quote}
\textbf{\pa{ਵੱਡੇ} \pa{ਕਮਰੇ} \pa{ਵਿੱਚ}} — \emph{vaḍḍe kamre vic} — in the big room
\end{quote}

This is one bounded location phrase. Notice \textbf{\pa{ਵੱਡੇ}} beside \textbf{\pa{ਕਮਰੇ}}; a full adjective-agreement table comes later, not here.
\end{grammarlens}

\subsection*{Your turn}
\begin{itemize}
\raggedright
  \item \emph{Hear:} the big-room phrase; YOU SAY: \textbf{\pa{ਵੱਡੇ} \pa{ਕਮਰੇ} \pa{ਵਿੱਚ}}
  \item \emph{Read it:} \textbf{\pa{ਵੱਡੇ} \pa{ਕਮਰੇ} \pa{ਵਿੱਚ}}
\end{itemize}

\begin{tcolorbox}[breakable,colback=teal!4,colframe=teal!35!black,title={Before you move on}]
Say \emph{in the big room}. (\textbf{\pa{ਵੱਡੇ} \pa{ਕਮਰੇ} \pa{ਵਿੱਚ}}.)

Source: UK Government, \href{https://www.hantsfire.gov.uk/wp-content/uploads/2021/03/Punjabi-Fire-Safety-in-the-Home.pdf}{Punjabi fire-safety guide}, uses \textbf{\pa{ਵੱਡੇ} \pa{ਘਰ} \pa{ਵਿੱਚ}}; Punjabi University, \href{https://punjabiuniversity.ac.in/Pages/Department.aspx?dsenc=129}{library description}, uses \textbf{\pa{ਕਮਰੇ} \pa{ਵਿੱਚ}} in its room phrase.
\end{tcolorbox}
\section[\texorpdfstring{\pa{ਉਸ} \pa{ਥਾਂ} \pa{ਦੇ} \pa{ਵੱਡੇ} \pa{ਕਮਰੇ} \pa{ਵਿੱਚ} (us thaan de …)}{ਉਸ ਥਾਂ ਦੇ ਵੱਡੇ ਕਮਰੇ ਵਿੱਚ (us thaan de …)}]{That place, its big room — one guided review}

\label{lesson:PA-C180-place-review}



\begin{quote}
Recall \emph{people}, \emph{that place}, then \emph{in the big room}: \textbf{\pa{ਲੋਕ}}, \textbf{\pa{ਉਸ} \pa{ਥਾਂ}}, \textbf{\pa{ਵੱਡੇ} \pa{ਕਮਰੇ} \pa{ਵਿੱਚ}}.
\end{quote}

\subsection*{Your turn: fit the known pieces}
\begin{quote}
\textbf{\pa{ਉਸ} \pa{ਥਾਂ} \pa{ਦੇ} \pa{ਵੱਡੇ} \pa{ਕਮਰੇ} \pa{ਵਿੱਚ}} — \emph{us thā̃ de vaḍḍe kamre vic} — in the big room of that place
\end{quote}

Use the whole phrase as a bounded model. \textbf{\pa{ਦੇ}} links the already known pieces here; this is not a new possessive paradigm.

\begin{itemize}
\raggedright
  \item \emph{Hear:} the place-and-room phrase; YOU SAY: \textbf{\pa{ਉਸ} \pa{ਥਾਂ} \pa{ਦੇ} \pa{ਵੱਡੇ} \pa{ਕਮਰੇ} \pa{ਵਿੱਚ}}
  \item \emph{Read it:} \textbf{\pa{ਉਸ} \pa{ਥਾਂ} \pa{ਦੇ} \pa{ਵੱਡੇ} \pa{ਕਮਰੇ} \pa{ਵਿੱਚ}}
  \item \emph{Say it:} \textbf{\pa{ਲੋਕ}}
\end{itemize}

\begin{tcolorbox}[breakable,colback=teal!4,colframe=teal!35!black,title={Before you move on}]
Rebuild the place-and-room phrase, then say who may be there.

Sources: Punjabi University, \href{https://dic.learnpunjabi.org/default.aspx?look=\%E0\%A8\%A5\%E0\%A8\%BE\%E0\%A8\%82}{RCPLT dictionary}, for \textbf{\pa{ਥਾਂ}}; Punjabi University, \href{https://punjabiuniversity.ac.in/Pages/Department.aspx?dsenc=129}{library description}, for \textbf{\pa{ਕਮਰੇ} \pa{ਵਿੱਚ}}.
\end{tcolorbox}
\section[\texorpdfstring{\pa{ਲੋਕ} \pa{ਉਸ} \pa{ਥਾਂ} \pa{ਦੇ} \pa{ਵੱਡੇ} … (lok us thaan de …)}{ਲੋਕ ਉਸ ਥਾਂ ਦੇ ਵੱਡੇ … (lok us thaan de …)}]{\pa{ਲੋਕ} \pa{ਉਸ} \pa{ਥਾਂ} \pa{ਦੇ} \pa{ਵੱਡੇ} \pa{ਕਮਰੇ} \pa{ਵਿੱਚ} \pa{ਹਨ।}}

\label{lesson:PA-C180-payoff}



\begin{quote}
Say \textbf{\pa{ਲੋਕ}}, then \textbf{\pa{ਉਸ} \pa{ਥਾਂ} \pa{ਦੇ} \pa{ਵੱਡੇ} \pa{ਕਮਰੇ} \pa{ਵਿੱਚ}}.
\end{quote}

\subsection*{Your turn: four-skill payoff}
\begin{quote}
\textbf{\pa{ਲੋਕ} \pa{ਉਸ} \pa{ਥਾਂ} \pa{ਦੇ} \pa{ਵੱਡੇ} \pa{ਕਮਰੇ} \pa{ਵਿੱਚ} \pa{ਹਨ।}} — people are in the big room at that place.
\end{quote}

\begin{itemize}
\raggedright
  \item \emph{Hear:} the one-sentence place notice; YOU SAY: \textbf{\pa{ਲੋਕ} \pa{ਉਸ} \pa{ਥਾਂ} \pa{ਦੇ} \pa{ਵੱਡੇ} \pa{ਕਮਰੇ} \pa{ਵਿੱਚ} \pa{ਹਨ।}}
  \item \emph{Read it:} \textbf{\pa{ਲੋਕ} \pa{ਉਸ} \pa{ਥਾਂ} \pa{ਦੇ} \pa{ਵੱਡੇ} \pa{ਕਮਰੇ} \pa{ਵਿੱਚ} \pa{ਹਨ।}}
  \item \emph{Write it:} \textbf{\pa{ਲੋਕ} \pa{ਉਸ} \pa{ਥਾਂ} \pa{ਦੇ} \pa{ਵੱਡੇ} \pa{ਕਮਰੇ} \pa{ਵਿੱਚ} \pa{ਹਨ।}} after covering the model
\end{itemize}

\begin{tcolorbox}[breakable,colback=teal!4,colframe=teal!35!black,title={Before you move on}]
Rebuild the line without looking. Keep the five forms together only in this bounded place note.

This is the third prerequisite tranche exposed by Punjabi A2 mock drafting. It does not unblock Punjabi A2 mock 1 by itself. It is not an A2 mock, not A2 readiness or attainment, and not external or human validation.

Sources: Punjabi University, \href{https://dic.learnpunjabi.org/default.aspx?look=\%E0\%A8\%B2\%E0\%A9\%8B\%E0\%A8\%95}{RCPLT dictionary}, for \textbf{\pa{ਲੋਕ}}; Punjabi University, \href{https://punjabiuniversity.ac.in/Pages/Department.aspx?dsenc=129}{library description}, for the public-room phrase; UK Government, \href{https://www.hantsfire.gov.uk/wp-content/uploads/2021/03/Punjabi-Fire-Safety-in-the-Home.pdf}{Punjabi fire-safety guide}, for \textbf{\pa{ਵੱਡੇ} ... \pa{ਵਿੱਚ}}.
\end{tcolorbox}
